\section{Relation to earlier work}\label{sec:related-work}

Schoenberg introduced P\'olya-frequency functions and connected translation
minors with bilateral Laplace transforms and variation-diminishing convolution
operators~\cite{schoenberg-ii,belton}. Karlin's monograph remains the
standard systematic account of total positivity and its composition
principles~\cite{karlin}; a modern review of PF functions and their preservers
is given by Belton, Guillot, Khare, and Putinar~\cite{belton}.

Finite-order translation kernels have a substantially wider range than the
infinite-order class. Khare's characterization of measurable \(\TN_p\)
functions for \(p\geq3\) provides a general finite-order setting~\cite{khare}.
The present proof instead derives a criterion specific to order four, using
successive quotient derivatives and a curvature transport identity.

The theta normalization used here is the classical one in de Bruijn's study of
the Riemann Fourier kernel~\cite{debruijn}. For the same kernel, Dimitrov and Xu
relate Wronskians of Fourier and Laplace transforms to correlation kernels and
real-zero criteria~\cite{dimitrov-xu}, while Csordas and Varga obtain strict
concavity and moment inequalities connected with the Riemann
Hypothesis~\cite{csordas-varga}. Grochenig explains the distinct Schoenberg
correspondence relevant to the Riemann Hypothesis: the PF function there is
recovered from the reciprocal of a Laguerre--P\'olya entire function, rather than
identified with the original Riemann kernel~\cite{grochenig}.

Micha\l{}owski first published an explicit negative \(5\times5\) Toeplitz
minor for this kernel, supported by an interval computation, together with
single-configuration positivity checks at orders at most
four~\cite{michalowski}. Global \(\PF_4\) is posed there as Open Problem 1.
Theorem~\ref{thm:main} resolves that problem and supplies an exact rational
negative order-five minor, completing the finite-order classification.
